\section{Split CPEL representation and the classical anti-diagonal}
\label{sec:cpel-classical-collapse}

The preceding recovery gates characterize when a four-valued evaluation returns to
\(\{\Tval,\Fval\}\).  Gate~8 asks a different question: what does that recovery look
like through the repository's existing split translation into the classical
probabilistic epistemic syntax \texttt{CPELFormula}?

The translation was already present syntactically.  Positive support and negative
support are sent to separate Boolean-tagged atomic channels,
\[
  p^+ \mapsto (p,\mathsf{true}),
  \qquad
  p^- \mapsto (p,\mathsf{false}),
\]
with negation exchanging the two translations, conjunction preserving positive
support conjunctively, and negative support translated through derived classical
disjunction.  Until this gate, however, the target syntax had no explicit evaluator
connected back to the finite 4-PEL model semantics.

\subsection{An explicit Boolean target semantics}

Gate~8 defines \texttt{evalCPEL} on \texttt{CPELFormula (Atom $\times$ Bool) Ag}.
At an atom, the Boolean tag selects either the positive or negative FDE support bit.
Classical negation and conjunction use Boolean negation and conjunction.  Belief uses
the same accessibility list, finite event-mass function, and threshold as the source
model, but thresholds the Boolean truth set of the translated subformula.

The derived CPEL disjunction then evaluates as ordinary Boolean disjunction.  This
small bridge is enough to state the first main result without any recovery assumption.

\begin{theorem}[Exact split-translation semantics]
\label{thm:cpel-split-bits}
For every legacy formula \(\varphi\), finite 4-PEL model \(m\), and world \(w\),
\[
  \llbracket \operatorname{tr}_{+}(\varphi)\rrbracket^{C}_{m,w}
    = \bigl(\llbracket\varphi\rrbracket_{m,w}\bigr)^{+},
\]
and
\[
  \llbracket \operatorname{tr}_{-}(\varphi)\rrbracket^{C}_{m,w}
    = \bigl(\llbracket\varphi\rrbracket_{m,w}\bigr)^{-}.
\]
\statusverified
\end{theorem}

The proof is by structural induction through atoms, negation, conjunction, and
Lockean belief.  The belief case is important: the induction hypotheses identify the
translated Boolean event lists pointwise with the source positive and negative
support-event lists, so the same threshold comparison is recovered on both sides.
No classicality, probability-integrity, or recovery premise is required.

Package the two target truth values into one FDE-shaped pair,
\[
  \operatorname{Pair}_{C}(m,w,\varphi)
  =
  \left(
    \llbracket \operatorname{tr}_{+}(\varphi)\rrbracket^{C}_{m,w},
    \llbracket \operatorname{tr}_{-}(\varphi)\rrbracket^{C}_{m,w}
  \right).
\]
The bitwise theorem immediately yields a pointwise representation theorem.

\begin{theorem}[Two-Boolean representation]
\label{thm:cpel-pair-representation}
For every \(m,w,\varphi\),
\[
  \boxed{
  \llbracket\varphi\rrbracket_{m,w}
  =
  \operatorname{Pair}_{C}(m,w,\varphi).
  }
\]
Thus the split target semantics represents all four FDE values, not only the
classically recovered ones.
\statusverified
\end{theorem}

This theorem should be read semantically and pointwise.  It does not identify the
source and target proof systems, and it is not a strong-completeness theorem.

\subsection{The classical anti-diagonal}

Because an FDE value is literally a pair of Boolean support coordinates, its four
vertices are
\[
  \Nval=(0,0),\qquad
  \Tval=(1,0),\qquad
  \Fval=(0,1),\qquad
  \Bval=(1,1).
\]
The classically recovered values are exactly the two vertices on which one coordinate
is the Boolean complement of the other.

\begin{theorem}[Classical anti-diagonal]
\label{thm:classical-antidiagonal}
For every FDE value \(v\),
\[
  v\in\{\Tval,\Fval\}
  \quad\Longleftrightarrow\quad
  v^{-}=\neg v^{+}.
\]
Consequently, for every formula,
\[
  \boxed{
  \llbracket\varphi\rrbracket_{m,w}\in\{\Tval,\Fval\}
  \quad\Longleftrightarrow\quad
  \llbracket\operatorname{tr}_{-}(\varphi)\rrbracket^{C}_{m,w}
  =
  \neg\llbracket\operatorname{tr}_{+}(\varphi)\rrbracket^{C}_{m,w}.
  }
\]
\statusverified
\end{theorem}

The first equivalence is independent of formula syntax.  The second combines it with
the exact split-translation theorem.  In the Boolean evidence square, the classical
sector is therefore the discrete anti-diagonal
\[
  \{(1,0),(0,1)\}
  =
  \{(p,n)\in\{0,1\}^{2}: n=\neg p\}.
\]
The two nonclassical vertices fail the same equation in opposite ways:
\(\Nval\) lacks both channels, while \(\Bval\) contains both.

This gives a precise sense in which classical recovery removes one independent
Boolean degree of freedom.  Outside the recovered slice, positive and negative
support are genuinely separate coordinates.  On the anti-diagonal, choosing the
positive coordinate fixes the negative coordinate.  The word ``degree'' here refers
to discrete Boolean freedom, not vector-space dimension.

\subsection{Gate~7 recovery collapses the split representation}

Gate~7 supplies a sufficient recursive condition for classicality under
\texttt{ModelProbabilityIntegrity}.  Combining that theorem with the anti-diagonal
characterization turns semantic recovery into an explicit collapse result.

\begin{corollary}[Probability-recovery split collapse]
\label{cor:probability-recovery-cpel-collapse}
Let \(m\) satisfy \texttt{ModelProbabilityIntegrity}, and let \(\varphi\) satisfy
\texttt{Formula.ProbabilityRecoveryAdmissible}.  Then at every world \(w\),
\[
  \llbracket\operatorname{tr}_{-}(\varphi)\rrbracket^{C}_{m,w}
  =
  \neg\llbracket\operatorname{tr}_{+}(\varphi)\rrbracket^{C}_{m,w}.
\]
Moreover, the full recovered FDE value can be reconstructed from the positive target
bit alone:
\[
  \llbracket\varphi\rrbracket_{m,w}
  =
  \operatorname{ofClassicalBool}
  \left(
    \llbracket\operatorname{tr}_{+}(\varphi)\rrbracket^{C}_{m,w}
  \right),
\]
where \(\operatorname{ofClassicalBool}(1)=\Tval\) and
\(\operatorname{ofClassicalBool}(0)=\Fval\).
\statusverified
\end{corollary}

The logical shape of Gates~5--8 can now be displayed as a compression chain:
\[
\begin{array}{c}
\text{regular evidence / recursive recovery}\[2pt]
\Downarrow\\[2pt]
\llbracket\varphi\rrbracket\in\{\Tval,\Fval\}\[2pt]
\Downarrow\\[2pt]
\llbracket\operatorname{tr}_{-}(\varphi)\rrbracket^{C}
=\neg\llbracket\operatorname{tr}_{+}(\varphi)\rrbracket^{C}\[2pt]
\Downarrow\\[2pt]
\text{one Boolean coordinate reconstructs the recovered value.}
\end{array}
\]

\paragraph{Gate boundary.}
Gate~8 proves an explicit Boolean semantics for the existing split CPEL target, an
unconditional pointwise representation of both FDE support bits, an exact
characterization of the classical sector as the complement locus, and the collapse
to one Boolean coordinate under Gate~7 probability recovery.  It does \emph{not}
formalize a CPEL Hilbert system, target validity, a canonical model, strong
completeness, or decidability.  The former source comment suggesting that the split
translation by itself guaranteed those results has therefore been narrowed to an
intended reduction strategy rather than promoted to a theorem.
